\subsection{System Architecture Overview}
\label{sec:arch-overview}

The proposed IDS is organised as a staged pipeline in which each component exposes a narrow data contract to the next, enabling the offline training and online inference paths to share identical preprocessing and feature-engineering logic. Figure~\ref{fig:pipeline} summarises the five logical stages.

\begin{figure}[h]
\centering
\begin{tikzpicture}[
    node distance=0.65cm and 0.85cm,
    every node/.style={font=\footnotesize},
    stage/.style={draw, rounded corners=2pt, minimum width=2.3cm, minimum height=0.9cm, align=center, thick},
    arrow/.style={-Stealth, thick}
]
    \node[stage, fill=blue!8]   (capture)  {Traffic\\Capture};
    \node[stage, fill=blue!8,  right=of capture]  (extract)  {Flow\\Extractor};
    \node[stage, fill=green!10, right=of extract]  (prep)     {Preprocessing\\\& Scaling};
    \node[stage, fill=orange!15, right=of prep]    (model)    {ML Inference\\(RF / XGB / DNN)};
    \node[stage, fill=red!10,  right=of model]    (alert)    {Alert /\\Log Sink};

    \draw[arrow] (capture) -- node[above=-1pt,font=\scriptsize]{pcap} (extract);
    \draw[arrow] (extract) -- node[above=-1pt,font=\scriptsize]{flow row} (prep);
    \draw[arrow] (prep)    -- node[above=-1pt,font=\scriptsize]{$\mathbf{x}\in\mathbb{R}^{d}$} (model);
    \draw[arrow] (model)   -- node[above=-1pt,font=\scriptsize]{$\hat{y}, p$} (alert);
\end{tikzpicture}
\caption{End-to-end IDS pipeline. Solid arrows indicate the online data path; the offline training path reuses the green and orange stages on a static corpus.}
\label{fig:pipeline}
\end{figure}

\subsubsection*{Stages.}
Each stage is a replaceable unit with a well-defined interface:

\begin{description}[leftmargin=1.2em,itemsep=2pt,topsep=2pt]
    \item[Traffic Capture.] Raw packets are mirrored from a switch SPAN port or tap and written to a rolling PCAP buffer by \texttt{tcpdump} / \texttt{tshark}. For the training corpus this stage is already materialised as the PCAP files shipped with CIC-IDS2017.
    \item[Flow Extractor.] \textsc{CICFlowMeter}~\cite{ref_cicflowmeter} groups packets into bidirectional flows keyed by the 5-tuple and emits one record per flow containing the 78 numeric features described in Section~\ref{sec:bg-dataset}. A flow is flushed when its idle timeout expires or an explicit FIN/RST is observed.
    \item[Preprocessing \& Scaling.] Rows with \texttt{NaN}/$\pm\infty$ are dropped, zero-variance columns are removed, and the remaining numeric columns are transformed using a \emph{frozen} \texttt{RobustScaler} whose median and IQR statistics are persisted from the training split (Section~\ref{sec:bg-preproc}). Categorical labels are encoded via a persisted \texttt{LabelEncoder} for training-time consumption.
    \item[ML Inference.] The scaled feature vector is handed to the active model (Random Forest, XGBoost, or DNN). The model emits a class index $\hat{y}$ and, when supported, a calibrated probability vector $p \in [0,1]^{C}$.
    \item[Alert / Log Sink.] Non-\texttt{BENIGN} predictions above the per-class confidence threshold (tuned on validation; see Section~\ref{sec:inference}) are forwarded to the SIEM or blocking engine; all predictions are written to a rolling structured log for offline auditing and model-drift analysis.
\end{description}

\subsubsection*{Offline vs.\ Online Modes.}
The same implementation artefact serves both modes. In the \emph{offline} mode the pipeline is fed a historical PCAP, the scaler and encoder are \emph{fitted} on the training split, and the learner's \texttt{fit} method consumes the cleaned feature matrix end-to-end. In the \emph{online} mode the scaler and encoder are loaded read-only from the training snapshot, the learner's \texttt{predict} method is invoked per-flow (or per micro-batch; Section~\ref{sec:inference}), and no component of the pipeline performs any stateful update. This symmetry is deliberate---it eliminates the most common failure mode in production ML systems, namely training/serving skew caused by subtly different feature-engineering code paths in the two phases.

\subsubsection*{Data Contracts.}
Between every two adjacent stages we define a typed contract that is enforced by a lightweight schema validator:

\begin{description}[leftmargin=1.2em,itemsep=2pt,topsep=2pt]
    \item[Extractor $\to$ Preprocessing:] a named-tuple of 78 numeric fields plus flow-identifier metadata (source/destination IP and port, protocol, timestamp).
    \item[Preprocessing $\to$ Inference:] a dense \texttt{float32} vector $\mathbf{x} \in \mathbb{R}^{d}$ with $d = 78 + e$, where $e$ is the number of engineered features introduced in Section~\ref{sec:feat-eng}, plus the carried-through flow-identifier metadata.
    \item[Inference $\to$ Alert Sink:] the flow-identifier metadata, the predicted label $\hat{y}$, and, when available, the top-1 confidence $p_{\hat{y}}$; optionally the full probability vector is persisted for audit.
\end{description}

This schema-first design means that every stage can be swapped---replacing \textsc{CICFlowMeter} with a bespoke eBPF extractor, or the Random Forest with a newer learner---without touching any upstream or downstream code.
